% ch7_c (书页 299-310 / p314-p325)
%--B-TAIL--
%JOIN%自旋 $1/2$。一般的规律是：自旋 $S$ 与偏振模式在转动下保持不变所对应的角 $\theta$ 之间的关系为 $S = 360^{\circ}/\theta$。引力场——它的波以光速传播——在量子理论中应当导致无质量粒子。注意到我们所描述的偏振模式在 $x$-$y$ 平面内转动 $180^{\circ}$ 下不变，我们预期相应的粒子——引力子（gravitons）——是自旋 2 的。探测这样的粒子还遥遥无期（就算我们永远无法直接探测到它们，也不会令人惊讶），但任何一门像样的引力量子理论都应当预言它们的存在。

% ---- 书页 299 / p314 ----（本页第一段已在上面 tail 中接续，正文自第二段起）

事实上，从自旋为 2 的引力子理论出发，再加上一些简单性质的要求，就为\emph{导出}广义相对论完整的 Einstein 方程提供了一条漂亮的途径。设想从对称张量 $h_{\mu\nu}$ 的拉格朗日量 (7.9) 出发，但现在设想这个场“其实真的是”一个在 Minkowski 时空中传播的物理场，而不是动力学度规上的微扰。（这个拉格朗日量不包含与物质的耦合，不过把它加上并不困难。）现在再追加一条要求：让 $h_{\mu\nu}$ 与它自身的能量-动量张量（下文讨论）耦合，同时也与物质的能量-动量张量耦合。这会在作用量中诱导出高阶非线性项，随之又诱导出阶数更高的“能量-动量”项。反复重复这一过程，就会引入一个无穷多项的级数，但这个级数可以求和成一个简单的表达式——也许是因为你已经知道答案——即 Einstein--Hilbert 作用量（可能还带有一些高阶项）。在这一过程中我们发现，物质与唯一的组合 $g_{\mu\nu} = \eta_{\mu\nu} + h_{\mu\nu}$ 耦合。换句话说，仅仅要求一个自旋为 2 的场与能量-动量张量耦合，我们最终就得到了广义相对论完全非线性的全部辉煌。背景度规 $\eta_{\mu\nu}$ 变得完全不可观测。当然，广义相对论的某些整体几何面貌被这套程序掩盖了；它归根结底只是为 Einstein 方程提供论证的另一种方式。

在我们谈论这些趣事之际，不妨指出：引力波的行为为弦论为何会给出引力的量子理论提供了一条线索。考虑一根弦圈的基本振动模式，如图 7.7 所示。一根弦圈有三个能量最低的模式：

% ---- 书页 300 / p315 ----（承接上页段末"……三个能量最低的模式："）

一个整体的大小振荡，外加两种独立的、振荡成椭圆的方式。它们给出三个无质量自由度：一个自旋为 0 的粒子（伸缩子，dilaton）和一个无质量的自旋为 2 的粒子（引力子）。注意弦的振荡与检验粒子在引力波影响下的运动之间有着明显的相似；这绝非偶然，这正是量子化的弦不可避免地给出引力的原因。（弦论最初是作为强相互作用理论被研究的，但各种模型都不可避免地预言一个多余的无质量自旋 2 粒子；人们最终意识到，如果把这个理论当作引力的量子理论来看待，这个缺陷反而可以变成优点。）那个多余而无用的自旋 0（标量）模式反映了这样的事实：弦论实际上预言的是一个标量-张量引力理论（如 4.8 节所讨论），而不是普通的广义相对论。由于自然界中并没有观察到这种无质量标量，必定要有某种机制在低能下赋予这个标量质量。

%FIG{7.7}: {一根弦圈的三个基本振动模式。整体的"呼吸"模式（最左）在转动下不变，给出一个自旋为 0 的粒子。另外两个模式与引力波的两个偏振相匹配，代表一个无质量自旋 2 粒子的两个态。}

\section{引力波的产生}

有了线性化真空方程的平面波解，剩下要讨论的就是源对引力辐射的产生。为此有必要考虑与物质耦合的 Einstein 方程 $G_{\mu\nu} = 8\pi G T_{\mu\nu}$。由于 $T_{\mu\nu}$ 不为零，度规微扰除了表示引力波的应变张量之外，还将包含非零的标量和矢量分量；我们不能假定解取横无迹形式 (7.90)。相反，我们将保留完整的微扰 $h_{\mu\nu}$，求出在远离源处产生的引力波，在那里我们随后可以施加横无迹规范。

即便有源存在，我们仍然可以引入一些方便的简化。首先定义无迹反转微扰（trace-reversed perturbation），
\begin{equation*}
\bar{h}_{\mu\nu} = h_{\mu\nu} - \frac{1}{2}\eta_{\mu\nu}h.
\tag{7.117}
\end{equation*}
“无迹反转微扰”这个名字是有道理的，因为
\begin{equation*}
\bar{h} = \eta^{\mu\nu}\bar{h}_{\mu\nu} = -h.
\tag{7.118}
\end{equation*}
显然，我们可以从无迹反转形式重构出原来的微扰，因此没有丢失任何信息。还要注意，如果我们在远离任何源的真空区域、并且可以取横无迹规范，那么无迹反转微扰将等于原来的微扰，
\begin{equation*}
\bar{h}^{\mathrm{TT}}_{\mu\nu} = h^{\mathrm{TT}}_{\mu\nu}.
\tag{7.119}
\end{equation*}
同时，我们仍然可以自由选择某种规范。在规范变换 (7.14) 下，无迹反转微扰按下式变换：
\begin{equation*}
\bar{h}_{\mu\nu} \rightarrow \bar{h}_{\mu\nu} + 2\partial_{(\mu}\xi_{\nu)} - \partial_\lambda\xi^{\lambda}\eta_{\mu\nu}.
\tag{7.120}
\end{equation*}

% ---- 书页 301 / p316 ----

于是，通过选取满足
\begin{equation*}
\Box\xi_\mu = -\partial_\lambda\bar{h}^{\lambda}{}_{\mu},
\tag{7.121}
\end{equation*}
的规范参数 $\xi_\mu$，我们就可以令
\begin{equation*}
\partial_\mu\bar{h}^{\mu\nu} = 0.
\tag{7.122}
\end{equation*}
这个条件称为\textbf{Lorenz 规范}，与电磁学中经常使用的类似条件 $\partial_\mu A^\mu = 0$ 相仿。\footnote{注意拼写。这个“规范”（gauge）源自 Ludwig Lorenz（1829--1891），而更有名的“变换”（transformation）则由 Hendrick Antoon Lorentz（1853--1928）发明。参见 J.~D. Jackson 和 L.~B. Okun，\emph{Rev. Mod. Phys.} 73, 663 (2001)。}注意，原来的微扰在这个规范下并不是横向的；确切地说，我们有
\begin{equation*}
\partial_\mu h^{\mu\nu} = \frac{1}{2}\partial^\nu h.
\tag{7.123}
\end{equation*}
把无迹反转微扰的定义代入我们的 Einstein 张量表达式 (7.8)，并利用 Lorenz 规范条件，便得到一个非常简洁的表达式
\begin{equation*}
G_{\mu\nu} = -\frac{1}{2}\Box\bar{h}_{\mu\nu}.
\tag{7.124}
\end{equation*}
用原来的微扰 $h_{\mu\nu}$ 写出的相应表达式则要稍微杂乱一些；这正是引入 $\bar{h}_{\mu\nu}$ 的原因。于是，这个规范下的线性化 Einstein 方程就是每个分量各自满足的一个波动方程，
\begin{equation*}
\Box\bar{h}_{\mu\nu} = -16\pi G T_{\mu\nu}.
\tag{7.125}
\end{equation*}
这类方程的解可以用 Green 函数求得，处理方式与电磁学中的类似问题完全相同。这里我们将循着 Wald (1984) 的路子回顾一下方法的梗概。

d'Alembert 算符 $\Box$ 的 Green 函数 $G(x^\sigma - y^\sigma)$，是波动方程在存在 $\delta$ 函数源时的解：
\begin{equation*}
\Box_x G(x^\sigma - y^\sigma) = \delta^{(4)}(x^\sigma - y^\sigma),
\tag{7.126}
\end{equation*}
其中 $\Box_x$ 表示对坐标 $x^\sigma$ 的 d'Alembert 算符。这类函数的用处在于，诸如 (7.125) 这样的方程的一般解可以写成
\begin{equation*}
\bar{h}_{\mu\nu}(x^\sigma) = -16\pi G\int G(x^\sigma - y^\sigma)T_{\mu\nu}(y^\sigma)\, d^{4}y,
\tag{7.127}
\end{equation*}
这可以立即得到验证。（注意，不需要任何 $\sqrt{-g}$ 因子，因为我们的背景就是简单的平直时空。）(7.126) 的解当然早就被算出来了，它们可以被视为“推迟的”（retarded）或“超前的”（advanced），取决于它们代表的是在时间中向前还是

% ---- 书页 302 / p317 ----（承接上页段末"……向前还是"）

向后传播的波。我们感兴趣的是推迟 Green 函数，它代表来自所考虑点之过去的信号的累积效应。它由下式给出
\begin{equation*}
G(x^\sigma - y^\sigma) = -\frac{1}{4\pi|\mathbf{x}-\mathbf{y}|}\,\delta\!\left[|\mathbf{x}-\mathbf{y}| - (x^{0} - y^{0})\right]\theta(x^{0} - y^{0}).
\tag{7.128}
\end{equation*}
这里我们已经用粗体表示空间矢量 $\mathbf{x} = (x^{1}, x^{2}, x^{3})$ 和 $\mathbf{y} = (y^{1}, y^{2}, y^{3})$，其模为 $|\mathbf{x}-\mathbf{y}| = [\delta_{ij}(x^{i}-y^{i})(x^{j}-y^{j})]^{1/2}$。$\theta$ 函数 $\theta(x^{0}-y^{0})$ 在 $x^{0} > y^{0}$ 时等于 1，否则为零。(7.128) 的推导会使我们离题太远，但在任何一本关于电动力学或物理中的偏微分方程的标准教科书中都可以找到。

把 (7.128) 代入 (7.127)，我们可以利用 $\delta$ 函数完成对 $y^{0}$ 的积分，得到
\begin{equation*}
\bar{h}_{\mu\nu}(t, \mathbf{x}) = 4G\int \frac{1}{|\mathbf{x}-\mathbf{y}|}T_{\mu\nu}(t - |\mathbf{x}-\mathbf{y}|, \mathbf{y})\, d^{3}y,
\tag{7.129}
\end{equation*}
其中 $t = x^{0}$。“推迟时间”（retarded time）一词指的正是量
\begin{equation*}
t_r = t - |\mathbf{x}-\mathbf{y}|.
\tag{7.130}
\end{equation*}
(7.129) 的含义应当是清楚的：引力场在 $(t, \mathbf{x})$ 处的扰动，是过去光锥上的点 $(t_r, \mathbf{x}-\mathbf{y})$ 处的能量与动量源的影响之总和，如图 7.8 所描绘。

%FIG{7.8}: {引力场在 $(t, x^i)$ 处的扰动用过去光锥上的事件来计算。}

让我们取这个一般解，考虑引力辐射由一个相当遥远的、由非相对论性物质组成的孤立源发出的情形；随着讨论的推进，这些近似会得到更精确的表述。首先我们需要建立几条傅里叶变换的约定，在处理振荡现象时它们总能让我们省力不少。给定一个时空函数 $\phi(t, \mathbf{x})$，我们关心的是它仅对时间的傅里叶变换（及其逆变换），

% ---- 书页 303 / p318 ----

\begin{gather*}
\widetilde{\phi}(\omega, \mathbf{x}) = \frac{1}{\sqrt{2\pi}}\int dt\, e^{i\omega t}\phi(t, \mathbf{x}), \\
\phi(t, \mathbf{x}) = \frac{1}{\sqrt{2\pi}}\int d\omega\, e^{-i\omega t}\widetilde{\phi}(\omega, \mathbf{x}).
\tag{7.131}
\end{gather*}
对度规微扰作变换，我们得到
\begin{align*}
\widetilde{\bar{h}}_{\mu\nu}(\omega, \mathbf{x}) &= \frac{1}{\sqrt{2\pi}}\int dt\, e^{i\omega t}\bar{h}_{\mu\nu}(t, \mathbf{x}) \\
&= \frac{4G}{\sqrt{2\pi}}\int dt\, d^{3}y\, e^{i\omega t}\frac{T_{\mu\nu}(t - |\mathbf{x}-\mathbf{y}|, \mathbf{y})}{|\mathbf{x}-\mathbf{y}|} \\
&= \frac{4G}{\sqrt{2\pi}}\int dt_r\, d^{3}y\, e^{i\omega t_r}e^{i\omega|\mathbf{x}-\mathbf{y}|}\frac{T_{\mu\nu}(t_r, \mathbf{y})}{|\mathbf{x}-\mathbf{y}|} \\
&= 4G\int d^{3}y\, e^{i\omega|\mathbf{x}-\mathbf{y}|}\frac{\widetilde{T}_{\mu\nu}(\omega, \mathbf{y})}{|\mathbf{x}-\mathbf{y}|}.
\tag{7.132}
\end{align*}
在这一串式子中，第一个等式就是傅里叶变换的定义，第二行来自解 (7.129)，第三行是从 $t$ 到 $t_r$ 的变量代换，第四行则又是傅里叶变换的定义。

现在我们作如下近似：源是孤立的、遥远的，并且缓慢运动。这意味着我们可以把源看作中心位于（空间）距离 $r$ 处，源的不同部分位于距离 $r + \delta r$ 处，且 $\delta r \ll r$，如图 7.9 所示。由于源缓慢运动，发射的辐射其频率 $\omega$ 大都足够低，使得 $\delta r \ll \omega^{-1}$。（本质上，光穿越源的速度远远快于源自身各部分的速度。）在这些近似下，$e^{i\omega|\mathbf{x}-\mathbf{y}|}/|\mathbf{x}-\mathbf{y}|$ 这一项可以用 $e^{i\omega r}/r$ 代替，并移到积分号外。这剩下
\begin{equation*}
\widetilde{\bar{h}}_{\mu\nu}(\omega, \mathbf{x}) = 4G\frac{e^{i\omega r}}{r}\int d^{3}y\, \widetilde{T}_{\mu\nu}(\omega, \mathbf{y}).
\tag{7.133}
\end{equation*}

%FIG{7.9}: {一个尺寸为 $\delta r$ 的源，位于距观者 $r$ 处。}

% ---- 书页 304 / p319 ----

事实上，我们无须计算 $\widetilde{\bar{h}}_{\mu\nu}(\omega, \mathbf{x})$ 的所有分量，因为 Lorenz 规范条件 $\partial_\mu\bar{h}^{\mu\nu}(t, \mathbf{x}) = 0$ 在傅里叶空间中意味着
\begin{equation*}
\widetilde{\bar{h}}^{0\nu} = \frac{i}{\omega}\partial_i\widetilde{\bar{h}}^{i\nu}.
\tag{7.134}
\end{equation*}
因此我们只需关心 $\widetilde{\bar{h}}_{\mu\nu}(\omega, \mathbf{x})$ 的类空分量，然后再由 (7.134) 恢复出 $\widetilde{\bar{h}}^{0\nu}$。要做的第一件事是令 $\nu = j$，从 $\widetilde{\bar{h}}^{ij}$ 求出 $\widetilde{\bar{h}}^{0j}$，然后再用它从 $\widetilde{\bar{h}}^{i0}$ 求出 $\widetilde{\bar{h}}^{00}$。于是，根据 (7.133)，我们要对 $\widetilde{T}_{\mu\nu}(\omega, \mathbf{y})$ 的类空分量取积分。我们从反向分部积分开始：
\begin{equation*}
\int d^{3}y\, \widetilde{T}^{ij}(\omega, \mathbf{y}) = \int \partial_k(y^{i}\widetilde{T}^{kj})\, d^{3}y - \int y^{i}(\partial_k\widetilde{T}^{kj})\, d^{3}y.
\tag{7.135}
\end{equation*}
第一项是一个面积分，由于源是孤立的，它将消失；第二项则可以通过 $\partial_\mu T^{\mu\nu} = 0$ 的傅里叶空间版本与 $\widetilde{T}^{0j}$ 联系起来：
\begin{equation*}
-\partial_k\widetilde{T}^{k\mu} = i\omega\widetilde{T}^{0\mu}.
\tag{7.136}
\end{equation*}
于是，
\begin{align*}
\int d^{3}y\, \widetilde{T}^{ij}(\omega, \mathbf{y}) &= i\omega\int y^{i}\widetilde{T}^{0j}\, d^{3}y \\
&= \frac{i\omega}{2}\int (y^{i}\widetilde{T}^{0j} + y^{j}\widetilde{T}^{0i})\, d^{3}y \\
&= \frac{i\omega}{2}\int \left[\partial_l(y^{i}y^{j}\widetilde{T}^{0l}) - y^{i}y^{j}(\partial_l\widetilde{T}^{0l})\right] d^{3}y \\
&= -\frac{\omega^{2}}{2}\int y^{i}y^{j}\widetilde{T}^{00}\, d^{3}y.
\tag{7.137}
\end{align*}
第二行的依据是，我们知道左边在 $i$ 和 $j$ 上是对称的；而第三行和第四行不过是把反向分部积分和 $T^{\mu\nu}$ 守恒重复了一遍。习惯上定义源能量密度的\textbf{四极矩张量（quadrupole moment tensor）}，
\begin{equation*}
I_{ij}(t) = \int y^{i}y^{j}T^{00}(t, \mathbf{y})\, d^{3}y,
\tag{7.138}
\end{equation*}
它是每个等时面上的一个常张量。四极矩张量的整体归一化是一个约定问题，绝非普遍统一，因此在比较不同文献时要当心。用四极矩的傅里叶变换来表示，我们的解具有紧凑的形式
\begin{equation*}
\widetilde{\bar{h}}_{ij}(\omega, \mathbf{x}) = -2G\omega^{2}\frac{e^{i\omega r}}{r}\widetilde{I}_{ij}(\omega).
\tag{7.139}
\end{equation*}

% ---- 书页 305 / p320 ----

我们可以把它变换回 $t$，得到\textbf{四极公式（quadrupole formula）}
\begin{equation*}
\boxed{\,\bar{h}_{ij}(t, \mathbf{x}) = \frac{2G}{r}\frac{d^{2}I_{ij}}{dt^{2}}(t_r),\,}
\tag{7.140}
\end{equation*}
其中和前面一样，$t_r = t - r$。

因此，由一个孤立非相对论性物体产生的引力波，正比于能量密度四极矩在观者过去光锥与源相交处的二阶导数。相比之下，电磁辐射的主要贡献来自电荷密度变化的\emph{偶极}矩。这个差别可以追溯到引力的普适性。偶极矩的变化对应于密度中心的运动——在电磁学中是电荷密度，在引力中则是能量密度。虽然没有什么能阻止一个物体的电荷中心振荡，但孤立系统质心的振荡违反动量守恒。（你可以让一个物体上下晃动，但你和地球也会朝相反方向稍稍晃动以作补偿。）四极矩量度系统的形状，一般比偶极矩小；由于这个原因，再加上物质与引力耦合微弱，引力辐射通常比电磁辐射弱得多。

一个特别有意思的情形是双星（两颗相互绕转的恒星）发出的引力辐射。为简单起见，考虑两颗质量为 $M$ 的恒星，在 $x^{1}$-$x^{2}$ 平面内的圆轨道上绕转，与它们共同质心的距离为 $R$，如图 7.10 所示。我们将用 Newton 近似处理恒星的运动，在这种近似下我们可以像 Kepler 那样讨论它们的轨道。刻画圆轨道最容易的方式，是把引力与向外的“离心力”相等：
\begin{equation*}
\frac{GM^{2}}{(2R)^{2}} = \frac{Mv^{2}}{R},
\tag{7.141}
\end{equation*}
由此得到
\begin{equation*}
v = \left(\frac{GM}{4R}\right)^{1/2}.
\tag{7.142}
\end{equation*}
完成一整圈轨道所需的时间很简单，
\begin{equation*}
T = \frac{2\pi R}{v},
\tag{7.143}
\end{equation*}
但对我们更有用的是轨道的角频率，
\begin{equation*}
\Omega = \frac{2\pi}{T} = \left(\frac{GM}{4R^{3}}\right)^{1/2}.
\tag{7.144}
\end{equation*}

% ---- 书页 306 / p321 ----

%FIG{7.10}: {一个双星系统。两颗质量为 $M$ 的恒星在 $x^1$-$x^2$ 平面内以轨道半径 $R$ 绕转。}

用 $\Omega$ 我们可以写出恒星 $a$ 的显式轨迹，
\begin{equation*}
x^{1}_{a} = R\cos\Omega t, \qquad x^{2}_{a} = R\sin\Omega t,
\tag{7.145}
\end{equation*}
以及恒星 $b$ 的轨迹，
\begin{equation*}
x^{1}_{b} = -R\cos\Omega t, \qquad x^{2}_{b} = -R\sin\Omega t.
\tag{7.146}
\end{equation*}
相应的能量密度为
\begin{equation*}
T^{00}(t, \mathbf{x}) = M\delta(x^{3})\left[\delta(x^{1} - R\cos\Omega t)\delta(x^{2} - R\sin\Omega t) + \delta(x^{1} + R\cos\Omega t)\delta(x^{2} + R\sin\Omega t)\right].
\tag{7.147}
\end{equation*}
这些密密麻麻的 $\delta$ 函数使我们能够直接积分，由 (7.138) 得到四极矩：
\begin{align*}
I_{11} &= 2MR^{2}\cos^{2}\Omega t = MR^{2}(1 + \cos 2\Omega t) \\
I_{22} &= 2MR^{2}\sin^{2}\Omega t = MR^{2}(1 - \cos 2\Omega t) \\
I_{12} &= I_{21} = 2MR^{2}(\cos\Omega t)(\sin\Omega t) = MR^{2}\sin 2\Omega t \\
I_{i3} &= 0.
\tag{7.148}
\end{align*}
由此再利用 (7.140)，容易得到度规微扰的分量：

% ---- 书页 307 / p322 ----

\begin{equation*}
\bar{h}_{ij}(t, \mathbf{x}) = \frac{8GM}{r}\Omega^{2}R^{2}
\begin{pmatrix}
-\cos 2\Omega t_r & -\sin 2\Omega t_r & 0 \\
-\sin 2\Omega t_r & \cos 2\Omega t_r & 0 \\
0 & 0 & 0
\end{pmatrix}.
\tag{7.149}
\end{equation*}
$\bar{h}_{\mu\nu}$ 的其余分量可以通过要求 Lorenz 规范条件得到满足而导出。

\section{引力辐射导致的能量损失}

到了这里，很自然想谈谈通过引力辐射发射的能量。然而这样的讨论立刻就会遇到种种问题，既有技术上的，也有哲学上的。正如我们前面提到过的，引力场中的能量没有真正的局域量度。当然，在弱场极限下——即我们把引力当作是在固定背景度规上传播的对称张量来描述——我们或许有望像对电磁学或其他任何场论那样，导出涨落 $h_{\mu\nu}$ 的能量-动量张量。在某种程度上这是可能的，但仍然困难。由于这些困难，文献中有许多不同的建议，关于在弱场极限下应当用什么作为引力的能量-动量张量；它们彼此不同，但对于诸如双星系统能量发射速率这类物理上适定的问题，它们给出的答案大体相同。

在技术层面上，当我们考虑能量-动量张量应当取什么形式时，困难开始出现。我们前面提到过电磁学和标量场论的能量-动量张量，两者共有一个重要特征——它们都是相关场的二次式。按照假设，我们处理弱场极限的办法是只保留度规微扰的线性项。因此，为了追踪引力波携带的能量，我们将不得不把计算扩展到至少 $h_{\mu\nu}$ 的二阶。事实上，我们一直都在稍稍作弊。在讨论引力波对检验粒子的效应以及双星系统产生波时，我们一直在利用检验粒子沿测地线运动这一事实。但我们知道，这是从能量-动量的协变守恒 $\nabla_\mu T^{\mu\nu} = 0$ 推出来的；然而在我们一直工作的阶次上，实际有的是 $\partial_\mu T^{\mu\nu} = 0$，而那将意味着检验粒子在平直背景度规中沿直线运动。这是弱场极限无法描述自引力系统的一个症状。在实践中，最好的办法是把弱场方程解到某个适当的阶，然后再事后论证解的有效性。我们将遵循 Misner, Thorne 和 Wheeler (1973) 第 35、36 章概述的程序，那里可以找到关于这些微妙之处的更多讨论。另见 Wald (1984) 与 Schutz (1985)。

现在让我们考察 Einstein 真空方程 $R_{\mu\nu} = 0$ 到二阶，看看结果如何能用引力场的能量-动量张量来解释。我们把度规和 Ricci 张量都展开，

% ---- 书页 308 / p323 ----

\begin{gather*}
g_{\mu\nu} = \eta_{\mu\nu} + h^{(1)}_{\mu\nu} + h^{(2)}_{\mu\nu} \\
R_{\mu\nu} = R^{(0)}_{\mu\nu} + R^{(1)}_{\mu\nu} + R^{(2)}_{\mu\nu},
\tag{7.150}
\end{gather*}
其中 $R^{(1)}_{\mu\nu}$ 取为与 $h^{(1)}_{\mu\nu}$ 同阶，而 $R^{(2)}_{\mu\nu}$ 与 $h^{(2)}_{\mu\nu}$ 则是 $(h^{(1)}_{\mu\nu})^{2}$ 阶的。因为我们在平直背景中工作，零阶方程 $R^{(0)}_{\mu\nu} = 0$ 自动满足。一阶真空方程就是
\begin{equation*}
R^{(1)}_{\mu\nu}[h^{(1)}] = 0,
\tag{7.151}
\end{equation*}
它确定一阶微扰 $h^{(1)}_{\mu\nu}$（相差规范变换）。二阶微扰 $h^{(2)}_{\mu\nu}$ 则将由二阶方程
\begin{equation*}
R^{(1)}_{\mu\nu}[h^{(2)}] + R^{(2)}_{\mu\nu}[h^{(1)}] = 0.
\tag{7.152}
\end{equation*}
确定。记号 $R^{(1)}_{\mu\nu}[h^{(2)}]$ 表示展开后的 Ricci 张量中与度规微扰呈线性关系的那些部分，也就是把由 (7.6) 给出的表达式应用于二阶微扰 $h^{(2)}_{\mu\nu}$；与此同时，$R^{(2)}_{\mu\nu}[h^{(1)}]$ 代表展开后 Ricci 张量的二次部分，
\begin{align*}
R^{(2)}_{\mu\nu} ={}& \frac{1}{2}h^{\rho\sigma}\partial_\mu\partial_\nu h_{\rho\sigma} + \frac{1}{4}(\partial_\mu h_{\rho\sigma})\partial_\nu h^{\rho\sigma} + (\partial^\sigma h^{\rho}{}_{\nu})\partial_{[\sigma}h_{\rho]\mu} - h^{\rho\sigma}\partial_\rho\partial_{(\mu}h_{\nu)\sigma} \\
&+ \frac{1}{2}\partial_\sigma(h^{\rho\sigma}\partial_\rho h_{\mu\nu}) - \frac{1}{4}(\partial_\rho h_{\mu\nu})\partial^\rho h - \left(\partial_\sigma h^{\rho\sigma} - \frac{1}{2}\partial^\sigma h\right)\partial_{(\mu}h_{\nu)\rho},
\tag{7.153}
\end{align*}
把它应用于一阶微扰 $h^{(1)}_{\mu\nu}$。没有交叉项，因为它们必然是更高阶的。

现在让我们把真空方程写成 $G_{\mu\nu} = 0$；这当然等价于 $R_{\mu\nu} = 0$，但能使我们以富有启发性的形式表达结果。在二阶我们有
\begin{equation*}
R^{(1)}_{\mu\nu}[h^{(2)}] - \frac{1}{2}\eta^{\rho\sigma}R^{(2)}_{\rho\sigma}[h^{(2)}]\eta_{\mu\nu} = 8\pi G t_{\mu\nu},
\tag{7.154}
\end{equation*}
其中我们定义了
\begin{equation*}
t_{\mu\nu} \equiv -\frac{1}{8\pi G}\left\{R^{(2)}_{\mu\nu}[h^{(1)}] - \frac{1}{2}\eta^{\rho\sigma}R^{(2)}_{\rho\sigma}[h^{(1)}]\eta_{\mu\nu}\right\}.
\tag{7.155}
\end{equation*}
关于这个表达式注意两点。第一，我们没有包含形如 $h^{(1)\rho\sigma}R^{(1)}_{\rho\sigma}[h^{(1)}]$ 的项，因为 $R^{(1)}_{\mu\nu}[h^{(1)}] = 0$。第二，(7.154) 的左边并不是完整的二阶 Einstein 张量，因为我们把涉及 $R^{(2)}_{\mu\nu}[h^{(1)}]$ 的项移到了右边，并大胆地把它们重新标记为一阶微扰的能量-动量张量 $t_{\mu\nu}$。这样一种认定看来非常有理：

% ---- 书页 309 / p324 ----（承接上页末"……这样一种认定看来非常有理："）

$t_{\mu\nu}$ 是一个对称张量，是 $h_{\mu\nu}$ 的二次式，它恰好以物质能量-动量张量会采取的方式表示微扰如何影响时空度规。（$h_{\mu\nu}$ 的线性项不起作用，因为 $G^{(1)}_{\mu\nu}[h^{(1)}]$ 已被一阶方程直接置为零。）注意，$t_{\mu\nu}$ 也是守恒的，在背景平直空间的意义下，
\begin{equation*}
\partial_\mu t^{\mu\nu} = 0,
\tag{7.156}
\end{equation*}
这从 Bianchi 恒等式 $\partial_\mu G^{\mu\nu} = 0$ 就可以知道。

遗憾的是，把 $t_{\mu\nu}$ 解释为能量-动量张量存在一些局限。当然，在完整的理论中它根本不是张量，但根据假设我们对这一点不予理会。更重要的是，它在规范变换（无穷小微分同胚）下不是不变的，你可以通过直接计算来验证这一点。绕开这个困难的一种办法，是把能量-动量张量在若干个波长上取平均，我们用尖括号 $\langle\cdots\rangle$ 表示这一操作。这一做法兼有哲学上和实用上的优点。从哲学的观点看，我们知道，正因为可以在任意一点选取 Riemann 法坐标，我们不可能定义一种纯粹局域的（在每一点恰好用该点的度规及其一阶导数来定义的）、可靠的引力能量-动量量度。然而，如果在若干个波长上取平均，我们就有希望捕捉到一个小区域内足够多的物理曲率，从而给出一个规范不变的量度。从实用的角度看，任何导数项（相对于导数的乘积）平均后都为零，
\begin{equation*}
\langle\partial_\mu(X)\rangle = 0.
\tag{7.157}
\end{equation*}
因此我们就有理由在平均括号下分部积分，
\begin{equation*}
\langle A(\partial_\mu B)\rangle = -\langle(\partial_\mu A)B\rangle,
\tag{7.158}
\end{equation*}
这将大大简化我们的表达式。

考虑到这些，让我们用二阶 Ricci 张量的表达式 (7.153) 来计算 (7.155) 定义的 $t_{\mu\nu}$。（从此以后我们将不再在度规微扰上写上标，因为我们只对一阶微扰感兴趣。）虽然取平均的部分动机是得到规范不变的答案，但实际计算相当繁难，所以为了演示，我们将在横无迹规范下进行，
\begin{equation*}
\partial^{\mu}h^{\mathrm{TT}}_{\mu\nu} = 0, \qquad h^{\mathrm{TT}} = 0.
\tag{7.159}
\end{equation*}
不要忘了，只有在真空中我们才被允许选取这个规范。在这个规范下，$R^{(2)\mathrm{TT}}_{\mu\nu}$ 的非零部分可以写成
\begin{align*}
R^{(2)\mathrm{TT}}_{\mu\nu} ={}& \frac{1}{2}h^{\rho\sigma}_{\mathrm{TT}}\partial_\mu\partial_\nu h^{\mathrm{TT}}_{\rho\sigma} + \frac{1}{4}(\partial_\mu h^{\mathrm{TT}}_{\rho\sigma})\partial_\nu h^{\rho\sigma}_{\mathrm{TT}} + \frac{1}{2}\eta^{\rho\lambda}(\partial^\sigma h^{\mathrm{TT}}_{\rho\nu})\partial_\sigma h^{\mathrm{TT}}_{\lambda\mu} \\
&- \frac{1}{2}(\partial^\sigma h^{\mathrm{TT}}_{\rho\nu})\partial^\rho h^{\mathrm{TT}}_{\sigma\mu} - h^{\rho\sigma}_{\mathrm{TT}}\partial_\sigma\partial_{(\mu}h^{\mathrm{TT}}_{\nu)\rho} + \frac{1}{2}h^{\rho\sigma}_{\mathrm{TT}}\partial_\sigma\partial_\rho h^{\mathrm{TT}}_{\mu\nu}.
\tag{7.160}
\end{align*}

% ---- 书页 310 / p325 ----

现在让我们施加平均括号，并在方便之处分部积分。(7.160) 的最后三项全部消失，因为分部积分给出的散度最终化为零。于是剩下
\begin{equation*}
\left\langle R^{(2)\mathrm{TT}}_{\mu\nu}\right\rangle = -\frac{1}{4}\left\langle(\partial_\mu h^{\mathrm{TT}}_{\rho\sigma})(\partial_\nu h^{\rho\sigma}_{\mathrm{TT}}) + 2\eta^{\rho\lambda}(\Box h^{\mathrm{TT}}_{\rho\nu})h^{\mathrm{TT}}_{\lambda\mu}\right\rangle.
\tag{7.161}
\end{equation*}
但微扰服从一阶运动方程，它令 $\Box h^{\mathrm{TT}}_{\mu\nu} = 0$。所以我们最终剩下
\begin{equation*}
\left\langle R^{(2)\mathrm{TT}}_{\mu\nu}\right\rangle = -\frac{1}{4}\left\langle(\partial_\mu h^{\mathrm{TT}}_{\rho\sigma})(\partial_\nu h^{\rho\sigma}_{\mathrm{TT}})\right\rangle.
\tag{7.162}
\end{equation*}
我们可以取迹得到标量曲率；分部积分之后我们再次遇到一个 $\Box h^{\mathrm{TT}}_{\mu\nu}$ 项，我们把它置为零，于是
\begin{equation*}
\left\langle\eta^{\mu\nu}R^{(2)\mathrm{TT}}_{\mu\nu}\right\rangle = 0.
\tag{7.163}
\end{equation*}
把这些表达式代入 (7.155)，就得到横无迹规范下引力波能量-动量张量的一个简单表达式：
\begin{equation*}
\boxed{\,t_{\mu\nu} = \frac{1}{32\pi G}\left\langle(\partial_\mu h^{\mathrm{TT}}_{\rho\sigma})(\partial_\nu h^{\rho\sigma}_{\mathrm{TT}})\right\rangle.\,}
\tag{7.164}
\end{equation*}
记住，在这个规范下非空间分量为零，$h^{\mathrm{TT}}_{0\nu} = 0$。因此你有时会看到上面的表达式用空间指标 $ij$ 而不是时空指标 $\rho\sigma$ 来写；两种写法显然等价。如果我们足够强大，能够不先选定规范就完成相应的计算，我们将会发现
\begin{align*}
t_{\mu\nu} = \frac{1}{32\pi G}\Big\langle(\partial_\mu h_{\rho\sigma})(\partial_\nu h^{\rho\sigma}) &- \frac{1}{2}(\partial_\mu h)(\partial_\nu h) \\
&- (\partial_\rho h^{\rho\sigma})(\partial_\mu h_{\nu\sigma}) - (\partial_\rho h^{\rho\sigma})(\partial_\nu h_{\mu\sigma})\Big\rangle.
\tag{7.165}
\end{align*}
稍作直接的处理便足以验证这个表达式其实是规范不变的，习题中会请你证明这一点。

让我们对单个平面波计算横无迹表达式 (7.164)，
\begin{equation*}
h^{\mathrm{TT}}_{\mu\nu} = C_{\mu\nu}\sin(k_\lambda x^\lambda).
\tag{7.166}
\end{equation*}
我们取了实部，并任意设定相位，使这个波是正弦而非余弦。于是能量-动量张量为
\begin{equation*}
t_{\mu\nu} = \frac{1}{32\pi G}k_\mu k_\nu C_{\rho\sigma}C^{\rho\sigma}\left\langle\cos^{2}(k_\lambda x^\lambda)\right\rangle.
\tag{7.167}
\end{equation*}

% 本块末段在原书书页 310 末 (7.167) 处截断，其后文字于书页 311 顶部继续；下一块应续接本段。
